\section{Background}

% \begin{figure*}[t]
% 	\subfloat[G1]{\includegraphics[width=0.24\linewidth]{flexheap/fig/back/g1_policy.pdf}}
% 	\hspace{0.5pt}
% 	\subfloat[VG1]{\includegraphics[width=0.35\linewidth]{flexheap/fig/back/vg1_policy.pdf}}
% 	\subfloat[FlexHeap]{\includegraphics[width=0.35\linewidth]{flexheap/fig/back/fh_policy.pdf}}
% 	\caption{Heap resizing policies of (a) G1, (b) VG1, and (c) \flexheapName{}.}
% 	\label{fig:back}
% \end{figure*}

This section provides background on the Garbage-First (G1) garbage collector, Vertical G1 (VG1) resizing mechanism, and the buffered I/O in Linux kernel.

\paragraph{\textbf{G1}}
It is the default garbage collector of OpenJDK designed to balance high application throughput with predictable pause times.
G1 divides the heap into equally sized regions and follows a generational model with two generations: young and old.
The young generation consists of eden and survivor regions.
Mutator threads allocate new objects in eden regions, and each young GC cycle reclaims eden space by evacuating live objects to survivor regions.
Objects that survive multiple young GCs are promoted to the old generation.

G1 performs both STW and concurrent GCs.
It begins with repeated young GCs (STW) that target only the young generation and pause all mutator threads.
As the old generation fills and its occupancy crosses a configurable threshold, G1 triggers a concurrent marking GC to identify live objects across the heap without stopping mutators.
After concurrent marking completes, G1 performs one or more mixed GCs (STW) that reclaim memory from both young and selected old regions with a high fraction of garbage.
Finally, full GCs (STW) are rare and serve as a fallback mechanism, pausing all mutators to reclaim the entire heap when incremental GCs cannot keep up with allocation demand.

G1 also relies on mutator write barriers to track cross-region references.
When a mutator updates an object reference, G1 marks the corresponding entry in a card table, a byte array in which each byte represents a \SI{512}{B} segment of the heap.
A marked card indicates that the corresponding segment may contain updated references that must later be examined by the garbage collector.
Concurrent refinement threads run alongside mutator threads and process these dirty cards while the application executes, identifying cross-region references and updating the remembered sets of the corresponding target regions.
By doing so, they move much of remembered-set maintenance out of the STW pause, thereby reducing GC pause time.
During young or mixed GC, when G1 collects a region, it uses that region’s remembered set to locate incoming cross-region references and update them to point to the objects’ new locations, without scanning the entire heap.

% \paragraph{\textbf{G1 resizing policy.}}
% At the end of each young, mixed, and full GC cycle, G1 performs heap resizing
% at the granularity of heap regions: shrinking involves uncommitting regions and
% returning them to the operating system, while growing reserves additional
% regions for future allocations. In the latest OpenJDK (version 26), heap
% resizing decisions follow different criteria depending on the type of
% collection (Figure~\ref{fig:back}(a)). For young and mixed GCs, resizing is
% based on the ratio of time spent in GC versus application execution, known as
% the GC time ratio (\texttt{-XX:GCTimeRatio}). If recent pause times are
% consistently low (e.g., the last eight consecutive GCs or the average of the
% last ten fall below a predefined threshold), G1 interprets this as low memory
% pressure and shrinks the heap. Conversely, if GC overhead is high (e.g., the
% last four GCs or the average of the last ten exceed the threshold), G1 grows
% the heap to improve application throughput. By contrast, for full GCs, resizing
% is based on occupancy. After a full collection, G1 evaluates the fraction of
% free space relative to runtime thresholds \texttt{-XX:MinHeapFreeRatio,
% 	-XX:MaxHeapFreeRatio} thresholds. If the free space falls below the minimum
% ratio, the heap is grown to provide more allocation headroom; if it exceeds the
% maximum ratio, the heap is shrunk by uncommitting regions. Thus, while young
% and mixed GCs pursue throughput goals, full GC resizing is driven by heap
% occupancy goal.
%
% In addition to resizing the heap, G1 also determines the size of the eden space
% (Figure~\ref{fig:back} (a)). This decision is key to balancing throughput and
% latency, as eden size directly impacts the number of allocations before the
% next GC. G1 uses a predictive model to estimate the pause time that would
% result from collecting a given number of eden regions, drawing on historical
% data such as evacuation cost per region, object copy time, and remembered-set
% processing overhead. If the estimated pause time exceeds the user-defined goal
% (\texttt{-XX:MaxGCPauseMillis}), G1 reduces the number of eden regions
% accordingly. This pause-time-driven sizing ensures that young GCs stay within
% latency targets while maximizing the interval between collections.

\paragraph{\textbf{VG1}}
It is a time-driven, occupancy-based mechanism that preserves G1's default heap resizing heuristics while returning memory to the OS more aggressively.
It triggers a concurrent mark cycle if no collection occurs within a user-configured interval (\texttt{-XX:G1PeriodicGCInterval}).
In G1, each concurrent marking cycle concludes with a remark phase, a short STW pause in which G1 finalizes the marking process to ensure that all live objects are correctly identified before mixed GCs.
At the end of the remark phase, VG1 makes heap resizing decisions based on heap occupancy.
If the gap between committed memory and the memory occupied by live objects exceeds a threshold, VG1 returns memory to the OS.

\paragraph{\textbf{Buffered I/O}}
Buffered I/O is the default file-I/O mode in most operating systems.
In Linux, it uses the kernel's I/O cache, which stores file data in memory in fixed-size pages (typically \SI{4}{KB}), to serve reads and temporarily buffer writes before they are flushed to storage.
This allows recently accessed file data to be served from DRAM instead of the storage device.
Buffered I/O also allows the OS to prefetch file pages into the I/O cache in anticipation of future accesses.
As a result, applications can issue file accesses at arbitrary offsets and sizes without managing low-level device requirements directly.
In contrast, direct I/O bypasses the kernel's I/O cache and transfers data directly between the application and the storage device.
Direct I/O typically requires applications to satisfy device-specific alignment and I/O-size constraints.
Consequently, applications using direct I/O do not benefit from additional I/O cache capacity.
